\section{Payment Protocol}
\label{sec:protocol}

\subsection{Identity, Certification, and Delivery}
\label{sec:protocol-certification}

Let $K_T$ be the encoded payment body and input claims, $\mathbf{C}_T$ the input commitments, and $\Sigma_T$ the output summary:
\begin{equation}
  \begin{aligned}
    K_T &= \operatorname{Encode}(T.\mathsf{Body},T.\mathsf{Claims}),\\
    \mathsf{id}(T) &= H(\texttt{TX\_V3},K_T,\mathbf{C}_T),\\
    c_T &= H(\texttt{DIRECT\_CERT\_V3},\operatorname{Encode}(\Sigma_T)).
  \end{aligned}
  \label{eq:protocol-identities}
\end{equation}
Owner signatures authenticate $\mathsf{id}(T)$ and stay outside its computation, and each output identifier binds this identity, the network, and the output position. Funding references and openings are encoded separately, so a restricted replacement leaves the authorized claims and outputs unchanged. Initial submission requires \texttt{OriginalFunding} for each claimed input with a valid opening, and a later representation preserves body, claims, commitments, and owner signatures. Verification caches therefore key complete command bytes.

Each payment carries
\[\begin{aligned}
\mathsf{IntentID}=\operatorname{Digest}(&\texttt{INTENT},\mathit{Network},\\
&\mathit{Subject},\mathit{Nonce}),
\end{aligned}\]
which names an operation of one payer on one network; it is not a nonce space shared among users. Members keep local intent records per organization, whereas public execution binds an intent to one TxID in the same atomic transition that executes the payment successfully. Two organizations can therefore certify, from independent inputs, two payments with the same subject and nonce and both obtain QCs, but at most one TxID executes publicly. This deduplicates transaction execution, not economic receipts from separately executed successors of conflicting QCs.

Under one Network and Subject, an honest wallet uses a fresh payment nonce for every new payment, across organizations and instances; a retry resubmits the same payment with its original authorization. A wallet must not change the nonce and reuse an old certificate, and two facts that share an intent are not merged into one guarantee. The experimental drivers derive nonces from input identities, which is not a general multi-device nonce allocation service.

A member checks ownership, routes, direct input certificates, value conservation, and the admission vector. One local transaction rechecks the input locks, final fee inputs, and the designated Worker's resource slices and records the approval with its exact debits; the approval keeps the transaction, admission vector, and direct input certificates of the request (Section~\ref{sec:protocol-resubmission}). CAL reservation equals $a_T$ including change even when every principal input is final, since each certified output may circulate before the payment executes. The member signs after the transaction commits, a retry reuses a non-invalidated approval, and releasing resource capacity keeps the consumption locks.

The collector sends one canonical request to four members and returns after verifying three signatures over the same $c_T$; different valid subsets give the same fact. The recipient checks the TXCer against the output body, index, owner, and trusted configuration, records the coin atomically, and can construct a successor request at once, carrying its direct input certificates in place of the ancestor chain.

Relay material is saved in the background, and bounded early tasks let \texttt{INSTALL} and public submission proceed independently. \texttt{INSTALL} stores complete material at members and records certified consumption locally; it creates neither a new approval nor a budget debit. Public submission proceeds separately. Members schedule backup submission relative to their first \texttt{INSTALL}, and repeated \texttt{INSTALL}s cannot postpone it. HTTP acceptance leaves the outbox entry pending until authenticated public completion removes it. Backup submission requires a retained complete submission; Section~\ref{sec:protocol-resubmission} adds a path from the approvals of the original signers.

The public \texttt{DirectSubmission} carries the transaction, admission vector, organization \texttt{SpendQC}, and direct input certificates. The committee reconstructs $\Sigma_T$ and verifies the same authorization; omitting a separate new-output TXCer wrapper changes the envelope and keeps the requirement. Public verification binds each input certificate to its network, issuer, fact, and output index. It requires exactly one certificate--index entry per certified input, rejects duplicate OutputIDs and extraneous entries, and permits distinct outputs of one parent certificate. Coverage is registered once per parent fact. The admission vector holds one CAL entry, bound to the issuer.
\subsection{Atomic Execution with Direct Obligations}
\label{sec:protocol-execution}

Public execution separates certificate-wide coverage from an actual missing-input obligation. On the first missing use of a parent fact $c$, \texttt{register} reserves its entire certified CAL amount and creates a Promise for every output. The coverage record maintains
\begin{equation}
  \mathsf{Original}_c=
  \mathsf{Discharged}_c+\mathsf{Paid}_c+\mathsf{Remaining}_c,
  \label{eq:protocol-coverage}
\end{equation}
where $\mathsf{Paid}_c$ is a gross cumulative counter and $\mathsf{Rec}_c\le\mathsf{Paid}_c$ counts the part repaid by source execution. Grant usage satisfies $\mathsf{Reserved}_g+\mathsf{Spent}_g\le B_g$, where $\mathsf{Spent}_g$ sums $\mathsf{Paid}_c-\mathsf{Rec}_c$ over the certificates under $g$. Here $\mathsf{Reserved}_g$ is public grant usage for registered coverage; members' local Reserved slices are signing reservations and may include unpublished certificates. Registration raises public grant usage and leaves the reserve account untouched; an effective compensation (Section~\ref{sec:protocol-decision}) debits the account and raises $\mathsf{Spent}_g$, and a repayment credits it and lowers $\mathsf{Spent}_g$ by the same amount. Only a missing output that is actually consumed opens a direct obligation and raises the public gap, so consuming one missing 40-CAL output of a 40/60 certificate reserves 100~CAL, opens a 40-CAL gap, and leaves the unused Promise covered.

Figure~\ref{fig:protocol-payment} gives the public transition. Its prestate includes earlier successful commands of the same block; static authorization may be cached, while consumption, balances, grants, and obligations are evaluated against the current state. All changes occur in a private overlay, and a failure discards the entire transition.

The intent check belongs to this evaluation: if the intent of $T$ is already bound to a different TxID, the transition returns a conflict. The binding is written in the same atomic transition that executes a payment successfully, and no production path deletes or rebinds it, so resubmitting the same bytes cannot remove the conflict.

\begin{figure}[t]
  \small
  \begin{algorithmic}[1]
    \REQUIRE Verified submission $T$, current state $S$, public time
    \STATE If $T$ is already settled, return no new effect
    \STATE $O\gets\operatorname{Overlay}(S)$
    \STATE Check intent, grants, and work; fund fee escrow in $O$
    \FOR{each CAL input $(x,i)$}
      \STATE Require $(x,i)$ to be unconsumed
      \IF{a final creation for $(x,i)$ exists}
        \STATE Check its body and evidence against the claim
      \ELSE
        \STATE Require $i=0$ and a matching direct TXCer $c$
        \STATE Register all of $c$'s coverage if unseen
        \STATE Create unique $\mathcal{O}_x$; add gap and Pending
      \ENDIF
      \STATE Record consumption of $(x,i)$ in $O$
    \ENDFOR
    \STATE Require equal total CAL input and output values
    \STATE Record $T$ settled; advance its fee stages
    \STATE Resolve each output: fulfill an Open obligation; repay the
    issuer for a Compensated one and mark it Recovered; otherwise fulfill
    any registered Promise. Create final instance~0 unless repaid
    \STATE Return all changes and results as one transition
  \end{algorithmic}
  \caption{Public payment execution. Any failed check rejects all
    provisional changes. Output resolution follows
    Section~\ref{sec:protocol-closure}.}
  \label{fig:protocol-payment}
\end{figure}

An obligation binds the consumed output, parent certificate fact, issuer, amount, consuming transaction, and historical input location, and public block time anchors its deadline. For an obligation created at height $h$, BeginBlock at $h+1$ fixes the deadline to that block's time plus the configured timeout, provided the obligation is still Open. The current chain and mixed-load runs use 30\,s. A public clock-tick command can advance block time during idle periods; deadline expiry authorizes a valid decision but does not itself execute compensation. The consuming payment records its Pending input obligations and fee progress. User-funded fees consume authenticated final FUEL, return excess input value as change, and escrow the signed maximum, which covers registration, settlement, closure, and one compensation charge per certified input; unique stage and obligation identities prevent duplicate charging, and unused escrow returns to the payer when the input obligations close. The child outputs are final even while some inputs remain obligations, so a further transaction consumes them under the ordinary final-input checks.
\subsection{Source Resubmission from Public Certificates}
\label{sec:protocol-resubmission}

The block that executes a payment carries the quorum certificates authenticating its certified inputs and their creating payments. Each member follows the block in its authenticated block transaction. For every carried certificate of its own organization, it loads the approval indexed by the certificate's fact, unless it has observed the source's public execution or already holds an outbox entry for that fact, and checks that the stored transaction and admission vector reproduce the fact. It then assembles the stored transaction, the stored direct input certificates, and the exposed QC, and applies ordinary payment verification. On success it writes an outbox entry in the same atomic local update, staggering the first attempt by member index; the existing relay submits the entry after commit. Payments of other organizations trigger the same path, and a stored record that fails these checks is corrupted state that stops the follower.

The job resubmits the payment that the organization certified, so it needs no new owner or member signature, creates no approval, keeps input locks unchanged, and reserves no capacity. An existing outbox entry keeps its retry time, and execution of the source removes the entry. The path needs an original signer that holds the approval, follows the chain, and reaches the committee, and it needs that no global intent conflict prevents the source from executing; a source still absent at its deadline falls to compensation. It gives no fixed-time guarantee. Section~\ref{sec:security-decision} shows that at least two honest signers hold the material whenever an obligation is public.
\subsection{Compensation Decisions}
\label{sec:protocol-decision}

After its deadline, an Open obligation closes as Compensated (the implementation's \texttt{Repaired} status) through a public \texttt{CompensationDecision} $\delta=(\mathit{net},x,h,j,k)$: the network, the consumed output, and the height, transaction index, and input index of the consuming payment. The fixed-length canonical encoding holds neither a submitter signature nor adaptation data. Authority comes from the public obligation, and committed state fixes the target coordinates, so every submitter forms the same bytes for a given obligation. Committee workers scan the due index and submit due decisions without requesting adaptation shares.

Figure~\ref{fig:protocol-decision} gives the transition. The executor compares the request's network, position, and input slot with the obligation before it reads any historical block, so a mismatching request is rejected without loading history. It then reads the canonical consuming payment at $(h,j)$, checks its identity, fact, input, and claimed amount, and requires \texttt{OriginalFunding} for $x$ with a valid opening and a same-width reserve replacement; it computes no replacement opening and contacts no adapter. One overlay carries every change: the reserve debit, the Compensated status, Paid, Spent, the gap, the consumer's Pending count, the compensation charge (from the consumer's escrow to rewards), and, when no input obligation remains, the consumer's fee closure and refund, together with the immutable decision record $\mathcal D_x$ (obligation snapshot and reserve-debit identity) and the pending-representation entry.

\begin{figure}[t]
  \small
  \begin{algorithmic}[1]
    \REQUIRE Decision $\delta=(\mathit{net},x,h,j,k)$, state $S$,
    execution height and block time
    \STATE Require $h$ to be older than the execution height
    \STATE If $\mathcal D_x$ exists, return no effect when it equals
    $\delta$; otherwise reject
    \STATE Require $\mathcal O_x$ Open with consumer position $(h,j,k)$
    and network $\mathit{net}$
    \STATE Require the deadline of $\mathcal O_x$ to be reached by block time
    \STATE Load the canonical consuming payment at $(h,j)$; check its
    identity, fact, input, and claimed amount; require
    \texttt{OriginalFunding} for $x$ with a valid opening and a
    same-width reserve reference
    \STATE $O\gets\operatorname{Overlay}(S)$
    \STATE In $O$: debit the issuer reserve; set $\mathcal O_x$ Compensated;
    move its Promise to Paid; update Spent, the gap, and the consumer's
    Pending count and fee stages
    \STATE In $O$: record $\mathcal D_x$ and the pending-representation
    entry
    \STATE Return the changes and the compensation result
  \end{algorithmic}
  \caption{Compensation decision. The transition reads and requests no
    adaptation material. Any failed check rejects all provisional
    changes.}
  \label{fig:protocol-decision}
\end{figure}

The result carries the single per-output effect, which names the parent and consumer facts, amount, input position, and debit identity, together with any fee outputs. A repeated identical decision returns an unapplied result with an empty change set, also after the source's repayment, and a decision that differs from the record is rejected; once the source has executed, no new decision for its obligation commits. Members apply the effect from the authenticated result together with their cursor and locate original approvals for the Paid and Pending updates, and wallets apply each fee output once. The compensation credits no principal to the recipient or to an executed successor.
\subsection{Source Closure, Repayment, and Cumulative Reconciliation}
\label{sec:protocol-closure}

A source transaction arrives through normal payment execution, including validation and consumption of its own inputs, and its outputs resolve within that transition according to their obligations. Arrival closes an Open obligation as Fulfilled: it fulfills the Promise, discharges the coverage, and lowers the consumer's Pending count. An output that no payment consumed fulfills any registered Promise and becomes final instance~0. A Compensated obligation means that a reserve debit funded the consumer, so arrival repays that debit instead of creating a second output: the amount returns to the issuer account that paid it, the obligation becomes Recovered, and the repayment counter and the grant's Spent change together, while the compensation decision stays as recorded. An already consumed instance~0 remains consumed in every branch; a repaid output has no final creation. The source's own inputs may open new obligations against their direct issuers within the same atomic transition.

Members reconstruct progress from authenticated consecutive blocks and results, and only an existing local approval supplies a debit to restore. For its CAL debit, member $m$ computes
\begin{equation}
  \begin{aligned}
    t_{m,T}&=
    \begin{cases}
      a_T-\bigl(\widehat{p}_{m,T}-\widehat{\mathsf{Rec}}_{m,T}\bigr),
        &\text{if settled},\\
      0,&\text{otherwise},
    \end{cases}\\
    \delta_{m,T}&=t_{m,T}-\mathsf{Applied}_{m,T},
  \end{aligned}
  \label{eq:protocol-release}
\end{equation}
where ``settled'' means that the member has observed its own payment execute successfully, and $\widehat{p}_{m,T}$ and $\widehat{\mathsf{Rec}}_{m,T}$ are the compensation and repayment recorded locally for that approval's outputs. On observing the execution of its payment, the member checks that the repayment in the authenticated result covers $\widehat{p}_{m,T}$ and sets $\widehat{\mathsf{Rec}}_{m,T}=\widehat{p}_{m,T}$, since that execution repays every compensated output of $T$; a member that approved $T$ after following the compensation recorded none. After checking bounds and monotonicity, it moves $\delta_{m,T}$ from Reserved to Available in the original Worker and updates Applied. Pending inputs leave CAL release unaffected, since their issuers carry those liabilities, and work resources wait for fee closure. An unexecuted certified source keeps its full CAL amount, including change, reserved at each approving member; compensation alone does not release it. Spent FUEL is not restored as unspent fee. Repeated blocks and late material cannot repeat a credit. A wallet records a repaid output as a terminal observation without certificate, final creation, or balance.

An authorized reserve increase debits an external, unprotected CAL account, credits the reserve, and raises its grant atomically. Each member adds the difference between the new and old $\lfloor2B_g/3\rfloor$ shares while Reserved and its approvals stay in place, and aggregate backing checks cover all grants that share an account.

\noindent\textbf{Superseded partial approvals.}
Equation~\eqref{eq:protocol-release} governs normal settlement and
repayment.  A separate local terminal transition applies when an
authenticated, successfully applied public payment under the same
organization configuration consumes a principal or user-funded fee
input instance that the member had locked for a different fact.
Before replacing the local spend candidate, the follower reloads the
immutable approval and checks the exact output instance, the distinct
fact and transaction, and the absence of observed execution, installed
or queued certificate material, pending obligations, compensation,
recovery, or any fee effect.

The member then returns, for every recorded local debit $d$, exactly
$d.\mathsf{Cap}-\mathsf{Applied}$ from Reserved to Available in the
original resource slice and records the approval as invalidated by the
first successful conflicting fact and height.  It does not set
\textsf{Settled}, close a fee escrow, modify the approval, clear other
input locks, or change public state.  The cumulative
$\mathsf{Applied}$ vector makes the transition idempotent when several
inputs of one approval conflict or a block is replayed. Retries and late
\texttt{INSTALL} are rejected. Exposure of an invalidated fact as a
purported certified source is treated as an accounting contradiction
and stops the follower under the stated fault model.
A timeout, \texttt{INSTALL} alone, or a global Intent
conflict between disjoint inputs does not authorize this release.
